\chapter{Conclusão}
\label{cap:conclusao}
A pesquisa investigou se a combinação de similaridade facial com similaridade visual global melhoraria a recuperação de fotografias de uma pessoa em comparação ao componente facial isolado. No protocolo Gallagher executado, com 20 consultas selecionadas e 588 candidatos por consulta, a resposta observada foi negativa para o ganho médio: a configuração facial atingiu mAP de 0,952261 e todas as fusões fixas avaliadas obtiveram valores inferiores. A fusão 0,9/0,1 apresentou benefícios em duas consultas, mas perdas em nove e empate em nove, resultando em mAP de 0,946274.

Essa conclusão é delimitada pelo desenho realizado. Ela não constitui demonstração de inutilidade do contexto, nem estabelece inferioridade universal de métodos de fusão. Mostra que a contribuição global, na forma de descritor ResNet50 e combinação linear sem calibração aprendida, não melhorou a média das consultas aceitas. O resultado contribui para evitar que a plausibilidade de uma hipótese seja confundida com sua confirmação empírica.

O objetivo de implementar um fluxo local de recuperação foi atendido por uma arquitetura que separa interface, operações de aplicação e adaptadores de extração e persistência. Foram integrados detecção e representação facial, descrição global, cálculo de escores e apresentação de fotografias ordenadas. A comparação controlada das configurações no Gallagher atendeu ao objetivo experimental principal, enquanto LFW e Holidays forneceram verificações auxiliares sob protocolos próprios. A análise pareada tornou explícitos os ganhos e as perdas por consulta que a média isolada não mostraria.

Em relação à proposta inicial, o trabalho realizou um recorte mais específico. Não foram concluídas uma comparação abrangente entre algoritmos de aprendizagem, uma implementação ativa de busca aproximada com FAISS ou uma avaliação controlada de escala e latência. A aceitação operacional não substitui essas etapas. A contribuição entregue consiste no protótipo, no protocolo de recuperação explicitado e na comparação rastreável de componentes sob condições registradas.

Entre as limitações centrais estão a quantidade reduzida e a dependência das consultas, a seleção condicionada à preparação bem-sucedida, a ausência de calibração independente da fusão e a impossibilidade de verificar separação completa entre dados de pré-treinamento e bases de avaliação. Os resultados auxiliares não usam, em todos os aspectos, os protocolos oficiais de suas bases. A persistência somente do topo dos rankings também restringe a auditoria independente de AP a partir dos artefatos atuais. Essas limitações não anulam a comparação realizada, mas restringem sua generalização.

Como continuidade científica, recomenda-se definir previamente novos conjuntos de consulta e validação, ampliar a diversidade de condições faciais e persistir rankings completos quando as condições de uso dos dados permitirem. Uma investigação de fusão deveria separar seleção de pesos e teste final, examinar a calibração dos escores e comparar pistas contextuais específicas. Estudos de desempenho computacional exigiriam desenho próprio, com controle de hardware, cache, repetições e tamanho de coleção. Uma avaliação de experiência de uso exigiria participação humana e os procedimentos institucionais correspondentes.

